\chapter{提案手法}\label{chap:method}

\section{定式化}

カテゴリ集合を $\mathcal{Y} = \{1, \dots, M\}$、符号語表を
$\bm{W} \in \{0,1\}^{M \times L}$ とする。入力 $\bm{x}$ に対する推定カテゴリは
\begin{equation}
  \hat{y} = \operatorname*{argmin}_{y \in \mathcal{Y}} \sum_{l=1}^{L} \bigl| w_{yl} - f_l(\bm{x}) \bigr|
  \label{eq:decode}
\end{equation}
で与えられる。

\begin{definition}
符号語表の任意の2行の間のハミング距離の最小値を、最小距離 $d_{\min}$ という。
\end{definition}

\begin{theorem}\label{thm:bound}
各2値分類器の誤り率が $p < 1/2$ で独立ならば、式\eqref{eq:decode}の誤り率は
$d_{\min}$ について単調に減少する。
\end{theorem}

\begin{proof}
2値分類器の誤り数が $\lfloor (d_{\min}-1)/2 \rfloor$ 以下なら正しく復号される。
この事象の確率は $d_{\min}$ について単調に増加する。
\end{proof}

\section{実験}

図\ref{fig:pipeline}に処理の流れを、表\ref{tab:result}に結果を示す。

\begin{figure}[tb]
  \centering
  \begin{tikzpicture}[box/.style={draw,rounded corners,minimum height=8mm,minimum width=22mm}]
    \node[box] (a) {入力 $\bm{x}$};
    \node[box,right=10mm of a] (b) {2値分類器 $f_l$};
    \node[box,right=10mm of b] (c) {復号};
    \draw[->] (a) -- (b);
    \draw[->] (b) -- (c);
  \end{tikzpicture}
  \caption{ECOC による分類の流れ}\label{fig:pipeline}
\end{figure}

\begin{table}[tb]
  \centering
  \caption{分類精度の比較}\label{tab:result}
  \begin{tabular}{lcc}
    \toprule
    手法 & 正解率 & 符号長 $L$ \\
    \midrule
    1-vs-Rest & 0.812 & 10 \\
    提案法    & \textbf{0.846} & 15 \\
    \bottomrule
  \end{tabular}
\end{table}
